\chapter{Step 07: An I/O Device with Buffers}
\label{ch:step07}

In step 06 every character that arrived on the serial line was sent back at
once, and then it was gone. A computer needs more: characters that arrive
while the program is busy must wait somewhere until the program asks for them,
and characters that the program writes must wait until the slow serial line
can take them. So we build an I/O device with an input buffer and an output
buffer. A buffer of this kind is a \term{FIFO}, a queue, and its content lives
in a memory block of the FPGA. On the way we meet three new pieces of
SystemVerilog: \cc{interface} and \cc{modport}, which bundle wires, and
\cc{package}, which bundles declarations. Everything is in the directory
\cc{07_io_interface}; the UART modules come from step 06
(Chapter~\ref{ch:step06}), the display from step 03.

\section{What the Design Does}
\label{sec:step07-design}

Figure~\ref{fig:step07-blocks} shows the parts. The module \cc{dev_io}
contains the UART receiver and transmitter and two FIFOs. Every character
that arrives on \cc{RX} goes two ways: into the input buffer, and into the
output buffer, from where it is sent back on \cc{TX}, so that you see what you
type. The other side of \cc{dev_io}, the one a computer would use, has two
operations:
\begin{description}
\item[getc] reads the oldest character of the input buffer and removes it
  (``pop'' it);
\item[putc] appends a character to the output buffer (``push'' it).
\end{description}
In this step there is no computer yet. The top module \cc{test} plays its
role with buttons: the display shows the oldest character in the input buffer
as a hexadecimal number, \cc{BTN1} pops it, and \cc{BTN2} pushes the byte
that is set at the inputs \cc{PMOD_1B} into the output buffer. \cc{LED1} is on
while the input buffer is empty, \cc{LED2} when it is full, and \cc{LED3}
while \cc{BTN1} is pressed although there is nothing to pop.

\begin{figure}[ht]
\centering
\begin{tikzpicture}[node distance=6mm and 9mm]
  \node[box,minimum width=17mm,minimum height=8mm] (rx) {\cc{uart_rx}};
  \node[box,fill=green!10,right=12mm of rx,minimum width=22mm,minimum height=8mm,yshift=9mm] (in) {input buffer\\[-2pt]\footnotesize FIFO, 512 bytes};
  \node[box,fill=green!10,right=12mm of rx,minimum width=22mm,minimum height=8mm,yshift=-9mm] (out) {output buffer\\[-2pt]\footnotesize FIFO, 512 bytes};
  \node[box,left=12mm of out,minimum width=17mm,minimum height=8mm,yshift=-12mm] (tx) {\cc{uart_tx}};
  \node[left=8mm of rx,note] (RX) {\cc{RX}};
  \node[left=8mm of tx,note] (TX) {\cc{TX}};
  \draw[arrow] (RX) -- (rx);
  \draw[arrow] (rx.east) -- ++(4mm,0) |- (in.west);
  \draw[arrow] (rx.east) -- ++(4mm,0) |- ([yshift=1.5mm]out.west);
  \draw[arrow] ([yshift=-1.5mm]out.west) -- ++(-3mm,0) |- (tx.east);
  \draw[arrow] (tx) -- (TX);
  \node[draw,dashed,black!50,rounded corners,fit=(rx)(in)(out)(tx),inner sep=3mm,
        label={[note]above:\cc{dev_io}}] (dev) {};
  \node[right=16mm of in,note,align=left] (getc) {getc:\\\cc{getc_pop},
    \cc{getc_en},\\\cc{getc_char}};
  \node[right=16mm of out,note,align=left] (putc) {putc:\\\cc{putc_push},
    \cc{putc_char}};
  \draw[arrow] (in) -- (getc);
  \draw[arrow] (putc) -- (out);
\end{tikzpicture}
\caption{The I/O device of step 07. The signals on the right are the side a
computer will use; in this step, buttons and the display take its place.}
\label{fig:step07-blocks}
\end{figure}

The new top module is short, because the work is done in \cc{dev_io}:

\svpart{firstline=14,lastline=30}{07_io_interface/test.sv}

Two things are new here: the line \cc{if_io if_io();} and the port
connection \cc{.io(if_io.server)}. They are the subject of the next section.

The ports of the display and of the input byte are vectors this time,
\cc{PMOD_1A} and \cc{PMOD_1B}, each of 8 bits. The pin file gives every bit its
own pin, with the index in brackets:
\begin{terminal}
set_io -nowarn PMOD_1A[0]   4
set_io -nowarn PMOD_1A[1]   2
...
set_io -nowarn PMOD_1B[7]  28
\end{terminal}
These are the same pins as \cc{A0} to \cc{A7} and \cc{B0} to \cc{B7}: the
display on PMOD 1A and the DIP switch on PMOD 1B, now addressed as two buses.

\section{Interfaces: Named Bundles of Wires}
\label{sec:step07-interface}

Between \cc{test} and \cc{dev_io} there are seven signals: three go into
\cc{dev_io}, four come out. We could list them as seven ports of \cc{dev_io}
and connect them one by one, and in step 08 we would do the same again in
the loader, and in step 09 in the processor. Instead we give the bundle a
name:

\svfile{07_io_interface/if_io.sv}

(The first two lines import a constant from a package; we come to them in the
next section. \cc{RAM_BYTE_SIZE} is 8.)

\textbf{An \cc{interface} is a named bundle of wires.} Between
\cc{interface if_io;} and \cc{endinterface} we declare signals, exactly as in
a module. The line \cc{if_io if_io();} in \cc{test} creates one
\term{instance} of the bundle, called \cc{if_io} as well: seven wires (or
groups of wires, \cc{getc_char} has eight) that exist once in the circuit.
Think of a ribbon cable with seven named strands. An interface contains no
logic of its own here; it is wiring.

\textbf{A \cc{modport} gives the direction of each wire, seen from one side.}
A ribbon cable has two ends. At the end that plugs into \cc{dev_io},
\cc{getc_pop} is an input and \cc{getc_char} an output; at the other end it is
the other way round. The modport \cc{server} lists the directions as
\cc{dev_io} sees them, the modport \cc{client} as the user of \cc{dev_io}
sees them (Figure~\ref{fig:step07-interface}). A module that takes the bundle
as a port names the modport:
\begin{lstlisting}
module dev_io #( ... ) (
    ...
    if_io.server io
);
\end{lstlisting}
Inside \cc{dev_io}, the strands of the bundle are called \cc{io.getc_pop},
\cc{io.getc_char}, and so on, and the modport tells the tools which of them
\cc{dev_io} may drive. When we instantiate \cc{dev_io}, we plug the
\cc{server} end of our bundle into this port with \cc{.io(if_io.server)}.
The top module \cc{test} uses the signals of the bundle directly, as
\cc{if_io.getc_pop} and so on; it takes the role of the client.

\begin{figure}[ht]
\centering
\begin{tikzpicture}[x=1mm,y=1mm]
  \node[box,minimum width=24mm,minimum height=44mm] (cl) at (0,0) {};
  \node[note,anchor=north] at (cl.north) {client side\\(here: \cc{test})};
  \node[box,minimum width=24mm,minimum height=44mm] (sv) at (88,0) {};
  \node[note,anchor=north] at (sv.north) {server side\\(\cc{dev_io})};
  \foreach \name/\y in {getc\_pop/9, putc\_push/4, putc\_char/-1} {
    \draw[arrow] (12,\y) -- (76,\y);
    \node[font=\scriptsize\ttfamily,fill=white,inner sep=1pt] at (44,\y) {\name};
  }
  \foreach \name/\y in {inbuf\_full/-6, getc\_en/-11, getc\_char/-16} {
    \draw[arrow] (76,\y) -- (12,\y);
    \node[font=\scriptsize\ttfamily,fill=white,inner sep=1pt] at (44,\y) {\name};
  }
  \draw[arrow] (76,-21) -- (12,-21);
  \node[font=\scriptsize\ttfamily,fill=white,inner sep=1pt] at (44,-21) {putc\_push\_done};
  \draw[decorate,decoration={brace,amplitude=3pt},thick,black!50] (18,12) -- (70,12)
    node[midway,above=2pt,note] {interface instance \cc{if_io}: seven wires};
  \node[note,anchor=north] at (0,-24) {\cc{modport client}:\\
    first three \cc{output}, rest \cc{input}};
  \node[note,anchor=north] at (88,-24) {\cc{modport server}:\\
    first three \cc{input}, rest \cc{output}};
\end{tikzpicture}
\caption{The interface \cc{if_io} as a bundle of wires between two modules.
The two modports describe the same wires, with the directions as seen from
either end.}
\label{fig:step07-interface}
\end{figure}

\textbf{What the tools make of it.} An interface is a convenience for us; in
the circuit there are only wires. sv2v translates interfaces into plain
Verilog by taking them apart. Here is an excerpt of \cc{test.v}:

\outputfile{step07_sv2v_if.txt}

The interface instance became a group of ordinary signals inside \cc{test}.
And every module that has an interface port, here \cc{dev_io} and \cc{fifo},
was removed as a module: sv2v copied its contents into the module that
instantiates it, with the interface signals connected directly. The circuit is
the same, but the hierarchy is gone. This is why \cc{make report} cannot list
\cc{dev_io} and \cc{fifo} separately (Section~\ref{sec:step07-resources}); the
report says so in its header.

\section{Packages and Include Files}
\label{sec:step07-package}

Several modules need the same constants, for example the width of a byte or,
from step 08 on, the number of address bits of the memory. We collect them in
a \term{package}:

\svfile{07_io_interface/pkg_ram.sv}

\textbf{A package is a named collection of declarations.} It describes no
hardware, it only gives names to types and constants. Inside it, the names are
declared with \cc{typedef enum}, which we know from step 06: there it named the
states of a state machine. Here the enumerations are used with explicit values
as a list of named integer constants: \cc{RAM_BYTE_SIZE} is 8,
\cc{RAM_QUAD_SIZE} is 64, \cc{RAM_ADDRW} is 17, and so on. The enumeration
\cc{data_type_t} and \cc{op_t} will be real types in step 08.

There are two ways to use a name from a package. The full name
\cc{pkg_ram::RAM_BYTE_SIZE} works everywhere. Or we \term{import} it once,
\cc{import pkg_ram::RAM_BYTE_SIZE;}, and then write just
\cc{RAM_BYTE_SIZE}. In \cc{if_io.sv} the import stands outside any module or
interface; for the tools we use, it then applies to everything that follows in
the same run.

\textbf{\cc{`include} copies text.} A package must be known before it is used.
The Makefile passes all \cc{*.sv} files to sv2v and Verilator in alphabetical
order, so \cc{if_io.sv} comes before \cc{pkg_ram.sv}. Therefore \cc{if_io.sv}
starts with \cc{`include "pkg_ram.sv"}. Like \texttt{\#include} in C, this is done
by a preprocessor before anything else: the line is replaced by the contents of
the file. Now the package is defined before \cc{if_io} uses it.

But it would be defined twice: once through the include, and once more when
the tools reach \cc{pkg_ram.sv} in the file list (and in step 08 many files
include it). A second definition of the same package is an error. The first
and the last two lines of \cc{pkg_ram.sv} prevent this; they are an
\term{include guard}. \cc{`ifndef PKG_RAM} means ``only if the macro
\cc{PKG_RAM} is not defined, keep the text up to \cc{`endif}''. The first copy
defines \cc{PKG_RAM}, so every later copy is dropped by the preprocessor. Just
as with interfaces, sv2v resolves all of this; \cc{test.v} contains no package
at all, only the constants substituted where they are used.

\section{A FIFO}
\label{sec:step07-fifo}

A FIFO (first in, first out) is a queue: we \term{push} elements at the back
and \term{pop} them at the front, and they come out in the order in which they
went in. In hardware we build it from a memory with $N$ entries and two
registers, the \term{read pointer} $r$ and the \term{write pointer} $w$. Push
stores the new element at address $w$ and increments $w$; pop takes the
element at address $r$ and increments $r$. When a pointer reaches the end of
the memory it starts again at address 0, so the memory is used as a ring
(Figure~\ref{fig:step07-ring}).

\begin{figure}[ht]
\centering
\begin{tikzpicture}
  \foreach \i in {0,...,7} {
    \pgfmathsetmacro{\a}{90-45*\i}
    \pgfmathsetmacro{\b}{90-45*(\i+1)}
    \ifnum\i>1 \ifnum\i<6
      \fill[green!15] (\a:1.1) arc (\a:\b:1.1) -- (\b:2.0) arc (\b:\a:2.0) -- cycle;
    \fi\fi
    \draw (\a:1.1) -- (\a:2.0);
    \pgfmathsetmacro{\m}{90-45*\i-22.5}
    \node[font=\scriptsize] at (\m:1.55) {\i};
  }
  \draw (0,0) circle (1.1); \draw (0,0) circle (2.0);
  \draw[arrow] (67.5-90-45:3.0) -- (67.5-90-45:2.1)
     node[pos=0,below right=-1mm,note] {read pointer $r=2$};
  \draw[arrow] (67.5-270:3.0) -- (67.5-270:2.1)
     node[pos=0,above left=-1mm,note] {write pointer $w=6$};
  \node[note] at (0,0) {$N=8$};
\end{tikzpicture}
\caption{A FIFO with $N=8$ entries as a ring. Entries 2 to 5 (shaded) hold
the queue; the next pop reads entry 2, the next push writes entry 6.}
\label{fig:step07-ring}
\end{figure}

\subsection*{Counting with pointers, modulo something}

How many elements are in the queue? If neither pointer had wrapped around, it
would be $w - r$. With wrapping, the pointers are numbers modulo $N$, and the
count is $(w - r) \bmod N$. Here is the catch: the count can be anything from
0 to $N$, which are $N+1$ values, but $(w - r) \bmod N$ has only $N$. An
empty FIFO has $w = r$; but after $N$ pushes into an empty FIFO, $w$ has gone
once around the ring, and again $w = r$. With pointers modulo $N$ we cannot
tell empty from full.

The usual solution, and the one used here, is to give each pointer one bit
more than the address needs. With $N = 2^k$ entries, an address has $k$ bits,
and the pointers have $k+1$ bits, so they count modulo $2N$. The memory
address is the pointer modulo $N$, which is simply its lower $k$ bits. The
number of elements is
\[
  n = (w - r) \bmod 2N ,
\]
and as long as $0 \le n \le N$ (we never push into a full FIFO and never pop
from an empty one), this is the true count: empty and full are now different
values of $w - r$. In terms of the bits:
\begin{itemize}
\item empty, $n = 0$: $w = r$, all $k+1$ bits equal;
\item full, $n = N$: $w - r \equiv N \pmod{2N}$, that is, the lower $k$ bits
  are equal (same address) and the top bits differ (one pointer is a lap
  ahead).
\end{itemize}

Our FIFO has $N = 512$ entries of one byte, so $k = 9$ and the pointers have
10 bits. An example: $r = 1000$ and $w = 6$ means $n = (6 - 1000) \bmod 1024
= 30$. The read pointer is at address $1000 \bmod 512 = 488$, the write
pointer at address 6: the queue occupies addresses 488 to 511 and 0 to 5.

\subsection*{The FIFO in SystemVerilog}

The FIFO has its own interface:

\svfile{07_io_interface/if_fifo.sv}

Like a module (Chapter~\ref{ch:step06}), an interface can have a parameter
list; here it fixes \cc{WIDTH = 8}. The FIFO itself:

\svfile{07_io_interface/fifo.sv}

The parameters compute the sizes: one memory block of $4096$ bits, entries
of $8$ bits, so $\mathtt{DEPTH} = 512$ entries and $\mathtt{ADDRW} =
\lceil \log_2 512 \rceil = 9$ address bits. \cc{$clog2} is evaluated by the
tools, not by the circuit: it computes a constant. All four are
\cc{localparam}s, which cannot be changed from outside: the FIFO always fills
exactly one memory block, and \cc{dp_bram4096} below computes the same depth
from the same width. The pointers are declared
as \cc{logic [ADDRW:0]}, that is with $\mathtt{ADDRW}+1 = 10$ bits, exactly as
in our analysis.

The \cc{always_ff} block describes the two pointer registers and their
next-state logic: each is a counter (register plus adder) whose increment is
enabled by \cc{push_back} or \cc{pop_front}. The \cc{case} over the
concatenation \cc{{fifo.push_back, fifo.pop_front}} lists all four
combinations; with both, both pointers advance and the count stays the same.
Wrapping around costs nothing: a 10-bit counter computes modulo $2^{10}$ by
itself, because the carry out of the top bit is dropped.

The \cc{always_comb} block computes \cc{full} and \cc{empty} with the
comparisons from above. \cc{read_ptr[ADDRW]} is the top bit, \cc{[ADDRW-1:0]}
the address bits.

Finally, the memory \cc{dp_bram4096} is instantiated with the write pointer's
address bits as write address and the read pointer's as read address. Its
write enable is \cc{push_back} itself: at the clock edge at which the pointer
advances, the byte is written into the memory.

The comment at the top states the rules for using the FIFO. The important one
for us: \cc{data_out} is the memory's output, and this output, as we see in a
moment, is a register that is loaded at every clock edge from the address that
the read pointer had \emph{before} the edge. So after a pointer changes, it
takes one more clock cycle until \cc{data_out} shows the new front element.

\section{Memory: Block RAM}
\label{sec:step07-bram}

A FIFO of 512 bytes needs $4096$ bits of storage. As flip-flops, that would
be $4096$ of our $5280$ logic cells, and the multiplexer that selects one of
512 bytes would need a few thousand LUTs more. This is what the memory blocks
mentioned in Section~\ref{sec:fpga} are for. The iCE40UP5K has 30
\term{block RAMs} (also called EBR, embedded block RAM) of $4$\,kbit each,
called \cc{SB_RAM40_4K} in the netlist. Each one can be configured as
$256 \times 16$, $512 \times 8$, $1024 \times 4$ or $2048 \times 2$ bits, and
it has two ports: one for writing and one for reading, each with its own
address, both working at the clock edge.

We do not instantiate \cc{SB_RAM40_4K} ourselves. Instead we describe a memory
in a way that matches what the block can do, and yosys recognises it:

\svfile{07_io_interface/dp_bram4096.sv}

\textbf{An array of vectors is a memory.} \cc{logic [WIDTH-1:0] mem [DEPTH];}
declares \cc{DEPTH} = 512 elements \cc{mem[0]} to \cc{mem[511]}, each a vector
of 8 bits. The dimension after the name, \cc{[DEPTH]}, is short for
\cc{[0:DEPTH-1]}. \cc{mem[addr_in]} selects an element by an address that is
a signal, so this is a memory with an address decoder, not a fixed wire.

\textbf{The pattern that becomes a block RAM.} In the \cc{always} block,
\begin{itemize}
\item one element is written at the clock edge, at one address, if an enable
  is 1: this is the write port;
\item \cc{data_out} is a register that takes, at the clock edge, the element at
  another address: this is the read port. The read is \emph{synchronous}: the
  data appears after the clock edge, not combinationally.
\end{itemize}
This is exactly the behaviour of \cc{SB_RAM40_4K}, whose read port has an
output register. (The block uses plain \cc{always} where we would write
\cc{always_ff}; it describes the same flip-flops.) yosys reports in
\cc{test.log} what it found:

\outputfile{step07_memlibmap.txt}

Both FIFO memories became block RAMs (\cc{$__ICE40_RAM4K_} is yosys's name for
\cc{SB_RAM40_4K} at this stage). The ROM of 16 entries that \cc{dev_hex}
contains, the table of digit patterns, is so small that yosys builds it from
logic cells instead. The statistics at the end of the log:

\outputfile{step07_yosys_stat.txt}

Two \cc{SB_RAM40_4K} for the two FIFOs. (The \cc{$scopeinfo} entries are not
hardware; yosys keeps them to remember the names of the parts of the design
after flattening. The flip-flops come in variants: \cc{E} has an enable
input, \cc{SR} and \cc{SS} a synchronous reset or set input. All of them are
the flip-flop of a logic cell.)

\textbf{Reading and writing the same address.} What should \cc{data_out} be
if, at one clock edge, we write address 5 and read address 5? With \cc{<=},
all right-hand sides are taken from before the edge (Chapter~\ref{ch:step02}),
so the description says: the old content (\term{read-first}). The block RAM
does not guarantee this; for such a collision its output is undefined. Left
to itself, yosys therefore adds logic around the block that detects a
collision and delivers the old value. The line before the declaration of
\cc{mem},
\begin{lstlisting}
(* no_rw_check *)
\end{lstlisting}
tells it not to. Text between \cc{(*} and \cc{*)} is an \term{attribute}: a
note for the tools, attached to the declaration that follows. It describes no
hardware itself; \cc{no_rw_check} promises yosys that we never rely on the
value read in a collision, so it may use the block RAM as it is. sv2v copies
attributes into \cc{test.v}. Whether the promise holds for our FIFOs is
Exercise~\ref{ex:step07-rwcheck}.

How much depends on the pattern? We synthesised \cc{dp_bram4096} alone, as
sv2v translates it, and in two modified versions: one without the attribute,
and one where the read is combinational, \cc{assign data_out = mem[addr_out];}:

\outputfile{step07_bram_variants.txt}

Without the attribute, 27 flip-flops and 19 LUTs appear next to the block
RAM: the collision logic (yosys's log calls it ``emulate read-first
behavior''), once per FIFO. The last version shows what happens if we leave
the pattern: the block RAM cannot read combinationally, so yosys builds the
memory from 4096 flip-flops and the read multiplexer from 4207 LUTs, more
than the whole chip has.

\begin{pitfall}
Synthesis directives are specific to a tool. An earlier version of
\cc{dp_bram4096.sv} had the comment
\cc{/* synthesis syn_ramstyle = "no_rw_check" */} after the declaration of
\cc{mem}. That is the syntax of another vendor's synthesis tool; yosys
ignores comments (and sv2v drops them), so the design silently carried the
46 extra cells per FIFO. If a directive matters, check in the statistics that
it has an effect.
\end{pitfall}

\textbf{Initial contents.} The \cc{initial} block sets \cc{mem[0]} to 255.
For a block RAM this is possible: its initial contents are part of the
bitstream, and all other entries start as 0. In our design this has a visible
effect: before the first character arrives, the read pointer is 0, the output
register shows \cc{mem[0]}, and the display shows \cc{FF}.

\section{The I/O Device}
\label{sec:step07-devio}

Now we can read \cc{dev_io}. It instantiates two interfaces of type
\cc{if_fifo} and two FIFOs, each connected through the modport \cc{fifo}:
\begin{lstlisting}
if_fifo inbuf(), outbuf();

fifo input_buffer(
    .clk(clk),
    .rst(rst),
    .fifo(inbuf.fifo)
);
\end{lstlisting}
Note that \cc{dev_io} has no wires of its own for the FIFO signals: \cc{inbuf}
\emph{is} the bundle of wires between \cc{dev_io} and the input buffer, and
\cc{dev_io} drives \cc{inbuf.push_back} and reads \cc{inbuf.data_out}.

\subsection*{From the receiver into the buffers}

\svpart{firstline=65,lastline=100}{07_io_interface/dev_io.sv}

The receiver from step 06 reports a byte by setting \cc{rx_ready} for one
clock cycle. The first block stores the byte in the register \cc{rx_data_buf}.
The second block remembers in the flip-flop \cc{rx_inbuf_push} that this byte
still has to go into the input buffer, and pushes it when the FIFO is not full
and has not been pushed in the previous cycle (rule (1) in the comment of
\cc{fifo.sv}). \cc{inbuf.push_back} is a register that is 1 for exactly one
cycle per byte. A third block of the same structure (lines 103--125) pushes
into the output buffer; besides the received bytes it accepts the bytes of
\cc{putc}, and it acknowledges them with \cc{putc_push_done}.

\begin{exercise}
\label{ex:step07-pointers}
In the input buffer, the read pointer is $r = 700$ and the write pointer is
$w = 188$ (both in decimal). How many characters are waiting? Is the buffer
full? At which memory addresses are the characters? What would happen if
\cc{dev_io} pushed 512 more characters without looking at \cc{full}?
\end{exercise}

\subsection*{When is a character available?}

\svpart{firstline=42,lastline=59}{07_io_interface/dev_io.sv}

Why not simply \cc{getc_en = !inbuf.empty}? Figure~\ref{fig:step07-timing}
shows the reason; it is taken from the simulation of this step, where a
character \cc{'Z'} (\cc{0x5a}) arrives. At the edge at time 14, the byte is
written into the memory and the write pointer advances, so \cc{empty} becomes
0. At the same edge, the output register of the memory loads the entry at the
read pointer, and since the write happens at the same edge, it gets the old
content, \cc{0xff}. Only at the next edge (time 16) does \cc{data_out} show
\cc{0x5a}. The register \cc{inbuf_empty_r} is \cc{empty} delayed by one
cycle, and \cc{getc_en} waits for both: it becomes 1 exactly when
\cc{getc_char} is valid.

The same delay occurs after a pop. In the clock cycle in which \cc{getc_pop}
is 1, \cc{getc_char} is the front character, and the pop takes effect at the
end of this cycle. In the cycle after it, the read pointer has moved, but the
memory output still shows the character just popped; only one edge later does
it show the next one. The register \cc{getc_pop_r}, \cc{getc_pop} delayed by
one cycle, covers this cycle: \cc{getc_en} is 0 in it, even if more
characters are waiting. So a user of \cc{dev_io} may rely on one simple rule:
whenever \cc{getc_en} is 1, \cc{getc_char} is the front character and a pop
is allowed. It takes the character in the cycle in which it pops (as the
comment in \cc{fifo.sv} says).

\begin{figure}[ht]
\centering
\begin{tikzpicture}[x=6.2mm,y=6.5mm]
  \foreach \name/\y in {CLK/8, rx\_ready/7, rx\_inbuf\_push/6, push\_back/5,
      write\_ptr/4, empty/3, data\_out/2, inbuf\_empty\_r/1, getc\_en/0} {
    \node[anchor=east,font=\small\ttfamily] at (5.7,\y+0.25) {\name};
    \draw[black!8,line width=3mm] (6,\y+0.25) -- (21,\y+0.25);
  }
  % clock: high at even times, rising edges at even times
  \draw[clocksig] (6,8.5) \foreach \t in {6,8,...,18} { -- (\t+1,8.5) -- (\t+1,8) -- (\t+2,8) -- (\t+2,8.5) } -- (21,8.5);
  \draw[signal] (6,7) -- (8,7) -- (8,7.5) -- (10,7.5) -- (10,7) -- (21,7);
  \draw[signal] (6,6) -- (10,6) -- (10,6.5) -- (12,6.5) -- (12,6) -- (21,6);
  \draw[signal] (6,5) -- (12,5) -- (12,5.5) -- (14,5.5) -- (14,5) -- (21,5);
  % write_ptr bus
  \draw[signal] (6,4) -- (13.8,4) -- (14.2,4.5) -- (21,4.5);
  \draw[signal] (6,4.5) -- (13.8,4.5) -- (14.2,4) -- (21,4);
  \node[font=\scriptsize] at (10,4.25) {0}; \node[font=\scriptsize] at (17.5,4.25) {1};
  \draw[signal] (6,3.5) -- (14,3.5) -- (14,3) -- (21,3);
  % data_out bus
  \draw[signal] (6,2) -- (15.8,2) -- (16.2,2.5) -- (21,2.5);
  \draw[signal] (6,2.5) -- (15.8,2.5) -- (16.2,2) -- (21,2);
  \node[font=\scriptsize] at (11,2.25) {ff}; \node[font=\scriptsize] at (18.5,2.25) {5a};
  \draw[signal] (6,1.5) -- (16,1.5) -- (16,1) -- (21,1);
  \draw[signal,draw=green!45!black] (6,0) -- (16,0) -- (16,0.5) -- (21,0.5);
  \foreach \t in {8,10,...,20} {
    \draw[black!12] (\t,-0.4) -- (\t,8.8);
    \node[below,font=\scriptsize] at (\t,-0.4) {\t};
  }
  \node[note,anchor=west] at (21.2,-0.7) {\cc{sim_time}};
\end{tikzpicture}
\caption{A character arrives in the input buffer (simulation of step 07; the
gray lines are rising clock edges). \cc{empty} becomes 0 one cycle before the
memory output shows the new character; \cc{getc_en} waits for it.}
\label{fig:step07-timing}
\end{figure}

\begin{pitfall}
A signal that says ``the data is valid'' must account for every reason why
the data can be stale. An earlier version of \cc{dev_io} had only
\cc{inbuf_empty_r}, so \cc{getc_en} was 1 in the cycle after a pop while
\cc{getc_char} still showed the popped character. On the board nobody
noticed: at 9600 baud a character takes more than 12\,000 clock cycles, and
the user had taken each one long before the next arrived. In a simulation in
which characters arrive every two clock cycles, the loader of
Chapter~\ref{ch:step08} read characters twice and skipped others.
\end{pitfall}

\subsection*{From the output buffer to the transmitter}

\svpart{firstline=143,lastline=159}{07_io_interface/dev_io.sv}

When the transmitter is idle and the output buffer is not empty, the block
sets \cc{pop_front} for one cycle. In this cycle \cc{data_out} is the front
byte, and at its end the byte is copied to \cc{tx_data_in} and \cc{tx_start} is
set: the transmitter from step 06 starts sending. Here, too, the byte is
taken in the cycle of the pop.

\subsection*{The top module}

\svpart{firstline=49,lastline=72}{07_io_interface/test.sv}

A button bounces (Chapter~\ref{ch:step02}): for a few milliseconds after it is
pressed or released, the contacts open and close several times. Each pop must
remove exactly one character, so \cc{test} does not look at \cc{BTN1}
directly. The module \cc{debounce} turns every press into a pulse
\cc{btn1_press} that is 1 for exactly one clock cycle, and in that cycle
\cc{getc_pop} is set to \cc{getc_en}: we pop only if there is something to
pop. Pushing with \cc{BTN2} works the same way, with a second instance of
\cc{debounce}. The \cc{initial} block sets the power-up value of the
flip-flop \cc{getc_pop}, which lives in the interface.

\svfile{07_io_interface/debounce.sv}

The debouncer is the one suggested in Chapter~\ref{ch:step04}, in three
parts:
\begin{itemize}
\item a \term{synchroniser} of two flip-flops, \cc{btn_s1} and \cc{btn_s2},
  because the button changes without regard to the clock;
\item the debounced \cc{level}, a flip-flop that follows \cc{btn_s2} only
  slowly. The counter \cc{cnt} counts the cycles in which \cc{btn_s2} differs
  from \cc{level}, and is cleared in every cycle in which they agree. Only if
  they differ in $2^W = 2^{17} = 131\,072$ cycles in a row, about 11\,ms at
  12\,MHz, does \cc{level} take the new value. A bounce shorter than that
  clears the counter again and has no effect. \cc{&cnt} is a \term{reduction}:
  the AND of all bits of \cc{cnt}, so it is 1 exactly when \cc{cnt} is
  $2^W - 1$, all ones.
\item an edge detector, \cc{level_r} one cycle behind \cc{level}, so
  \cc{press} is 1 in the first cycle in which \cc{level} is 1.
\end{itemize}
The price is a delay of 11\,ms between pressing the button and the pop, which
nobody notices, and about 60 cells per button
(Section~\ref{sec:step07-resources}).

\begin{exercise}
\label{ex:step07-bounce}
In \cc{test.sv}, instantiate both debouncers with \texttt{debounce \#(.W(2))}, so
that a level counts after four stable cycles. The testbench of the next
section presses each button once, but lets it bounce: four short pulses of
300 cycles before it stays pressed, and three short interruptions when it is
released. Predict what the testbench prints, and check with \cc{make sim}.
\end{exercise}

\section{Simulation}
\label{sec:step07-sim}

At 9600 baud a single bit lasts $12\,000\,000 / 9600 = 1250$ clock cycles, and
a character more than ten times as long. A simulation that sends real
characters through \cc{uart_rx} would mostly show waiting. The testbench
therefore bypasses the receiver and sets its outputs directly. It drives the
clock in a function \cc{tick}, one cycle per call, and at three rising edges
it delivers the characters \cc{'Z'}, \cc{'1'} and \cc{'2'}:

\cpppart{firstline=63,lastline=86}{07_io_interface/tb_test.cpp}

\cc{dut->test->io0->uart_rx0->rx_ready} is the signal \cc{rx_ready} inside
the instance \cc{uart_rx0} inside \cc{io0} inside \cc{test}. Verilator makes
such internal signals accessible from C++ only if they are marked with the
comment \cc{/* verilator public */}, which is why the copies of the UART
modules in this directory differ from those of step 06 in exactly this point.
In the same way, \cc{test.sv} marks \cc{hex_val}, the front character that the
display shows, so that the testbench can read it.

The rest of the testbench plays the user. It presses \cc{BTN1} once to pop
\cc{'Z'}, and then \cc{BTN2} once to push the byte \cc{0x41}, \cc{'A'}, set
at \cc{PMOD_1B}. To test the debouncers, each press bounces:

\cpppart{firstline=101,lastline=122}{07_io_interface/tb_test.cpp}

\cc{HOLD} is 140\,000 cycles, a little more than the $2^{17}$ that the
debouncer waits. Meanwhile, a small decoder in \cc{tick} samples \cc{TX} in
the middle of each bit and collects the bytes that the transmitter sends: the
echo of the three received characters, and the pushed \cc{'A'}. At the end the
testbench prints what it saw and compares it with what it expects:

\outputfile{step07_sim.txt}

One press, one pop, one push. The whole run is about 600\,000 clock cycles;
only the first 100 go into \cc{waveform.vcd}, which keeps the file small.
Figure~\ref{fig:step07-timing} is the beginning of it.

\section{Resources}
\label{sec:step07-resources}

\outputfile{step07_report.txt}

The design fits easily: 285 logic cells and 2 of the 30 block RAMs. The list
of modules shows the effect of sv2v from
Section~\ref{sec:step07-interface}: \cc{dev_io} and both FIFOs have interface
ports and were merged into \cc{test}, so their cells are counted there. The
two memories appear once as \cc{dp_bram4096}, with the single block RAM cell
we saw above, and the two debouncers once as \cc{debounce} with 62 cells (the
list shows each module once, not once per instance).

\begin{exercise}
\label{ex:step07-rwcheck}
The attribute \cc{(* no_rw_check *)} promises yosys that a read and a write
of the same address at the same edge never matter. Argue that the promise
holds for the I/O device: when can a write and a read of the same address
happen at the same edge, and is \cc{data_out} used in the following cycle?
Then delete the attribute and check with \cc{make report} what it saves.
\end{exercise}

\section{What Is Next}

We have a device that buffers input and output, and a defined way to read and
write characters through \cc{if_io}. What is missing for a computer is memory
for the program and its data: 512 bytes of block RAM would not get us far. The
next step uses the much larger memory blocks of the iCE40UP5K and writes a
program into them, typed in as hexadecimal digits through \cc{dev_io}.

\solutions

\solution
$n = (188 - 700) \bmod 1024 = 512$, so 512 characters are waiting and the
buffer is full: the address bits are equal ($700 \bmod 512 = 188$) and the
top bits differ ($700 \ge 512$, $188 < 512$). The characters occupy all
addresses, starting at the front at address 188, up to 511, then 0 to 187.
Pushing into a full buffer is forbidden by the rules of the FIFO, and
\cc{dev_io} checks \cc{full} before every push. If 512 pushes happened anyway,
they would overwrite all waiting characters, and $w$ would become
$188 + 512 = 700 = r$: the FIFO would report \cc{empty}, although 512
characters were pushed since the last pop.

\solution
With $W = 2$, a level counts after $2^2 = 4$ stable cycles. The bounces last
200 and 300 cycles, far more than four, so each of them is a separate press
for the debouncer: four presses before the button stays down, the real one,
and three more at the release, eight rising edges of \cc{level} per button.
The eight pops of \cc{BTN1} remove all three characters (pops on an empty
buffer are suppressed by \cc{getc_en}), and the display shows \cc{00}, the
never-written entry 3 of the memory. The eight pushes of \cc{BTN2} send eight
\cc{'A'}:

\outputfile{step07_sim_w2.txt}

The debouncer must wait longer than the contacts bounce; with $W = 17$ it
waits about 11\,ms.

\solution
A write to the address at the read pointer happens only when the FIFO is empty
(otherwise the write pointer is ahead of the read pointer, and when it is a
whole lap ahead, the FIFO is full and no push is allowed). After this edge,
\cc{data_out} may be wrong for one cycle. For the input buffer, \cc{getc_en}
stays 0 in this cycle because \cc{inbuf_empty_r} is still 1. For the output
buffer, the pop is decided at the end of this cycle and happens in the next,
when \cc{data_out} has been loaded again. So the device works without the
extra logic. Without the attribute, \cc{make report} shows 47 cells for
\cc{dp_bram4096} instead of 1: the attribute saves 46 cells per FIFO.
